\section{Conclusion}
\label{sec:conclusion}

The simulation results indicate that standard IEC/IEEE~60255-118-1 evaluation
tests may not fully characterize frequency estimator behavior in converter-dominated
grids, where phase discontinuities, high RoCoF, and harmonic distortion can
expose failure modes that isolated single-disturbance tests do not reveal. This
benchmark evaluates twelve estimators across five algorithmic families
under five stress scenarios and reports three main observations.

\begin{enumerate}

\item \textbf{Latency--Robustness Trade-off:}
Window-based methods achieve strong steady-state accuracy, but their
behavior differs across scenario types. IpDFT is better suited to
metering-oriented operation, whereas TFT reduces observation delay at
the cost of higher noise sensitivity under composite disturbances.
Loop-based estimators (SRF-PLL, SOGI-FLL) are computationally efficient
and handle isolated phase jumps well, but exhibit elevated trip-risk
under the composite Multi-Event sequence
($T_{\mathrm{trip}}=0.426$\,s for SRF-PLL).

\item \textbf{Model-Based Estimation:}
Among the model-based estimators, the \textbf{RA-EKF}---combining an
explicit $\dot{\omega}$ state with Innovation-Driven Covariance Scaling
and Event Gating---achieves a favorable trade-off between tracking
accuracy and trip-risk. Under Composite Islanding, the RA-EKF incurs
zero threshold breaches across all $N\!=\!30$ Monte Carlo trials
($T_{\mathrm{trip}}=0$), versus $T_{\mathrm{trip}}=0.120$\,s for the
standard EKF, at a computational cost of $68.9$\,$\mu$s/sample (Python
baseline; C/HIL validation remains future work).

\item \textbf{Limits of Standard Compliance Tests:}
Several estimators passing IEC/IEEE~60255-118-1 step and ramp tests
exhibited degraded performance under composite multi-disturbance
scenarios, suggesting that standard compliance tests may not fully
capture estimator resilience in IBR-dominated conditions.

\end{enumerate}

Data-driven approaches can be competitive in some scenarios, but
computational cost varies considerably. PI-GRU requires approximately
$145$\,ms/sample (Python/PyTorch CPU baseline), which is prohibitive for
embedded real-time protection without hardware acceleration. In the
evaluated scenarios, the pre-trained PI-GRU model did not outperform
scenario-optimized classical methods on RMSE or trip-risk, suggesting
that a fixed pre-trained model cannot substitute for scenario-specific
tuning under these disturbance conditions. This benchmark is entirely simulation-based; hardware-in-the-loop
and physical relay validation remain as next steps.

The results suggest that the IEC/IEEE~60255-118-1 test scope may not
fully characterize estimator resilience in IBR-dominated grids, as it
lacks composite phase-step evaluations under simultaneous harmonic
distortion. A \emph{Composite Phase-Step with Harmonic Distortion}
test---analogous to Scenarios~D--E---evaluated via probabilistic
compliance bounds ($\mathrm{PCB} = \mu_{|e|} + 3\sigma_{|e|}$) may
provide a more representative assessment of relay security under
realistic IBR conditions.